\section*{Präambel}

In den letzten Jahren wurden im deutschsprachigen Raum rege Debatten über die Qualität der Lehre geführt. Bei diesen war der Einfluß praxisferner Experten unserer Meinung nach überproportional groß. Dadurch fehlten Standpunkte und Erfahrungen derer, die die Praxis erleben, nämlich unsere. Als Studierende aller Semester, aller Hochschularten, aller mathematischen Studiengänge und aller auslaufenden oder eingeführten Studiensysteme sehen wir unsere Standpunkte im Ergebnis der Debatte nur unzureichend berücksichtigt.

Gleichzeitig muss man feststellen: Die Umsetzung des Reformprozesses ist hinter den Erwartungen der Hochschulen, Wirtschaft und Politik zurückgeblieben. Die gestiegene Belastung hindert die Studierenden an der Entfaltung ihrer wissenschaftlichen Fähigkeiten, worunter die gesamte Wissenschaftslandschaft leidet und nachhaltig negativ beinflusst wird.

Getrieben von einem Wunsch nach Eliten werden die Anforderungen der Masse vernachlässigt. Insbesondere wird ignoriert oder übersehen, dass es das Ziel der Hochschulen sein muss, nicht nur eine Elite zu schaffen, sondern das Gros erfolgreich zu Wissenschaftlern auszubilden. Hierzu müssen gewisse Anforderungen von jeder Hochschule erfüllt sein, eben die im Folgenden aufgeführten Minimalstandards. Die Bringschuld zur Erfüllung dieser Minimalstandards liegt sowohl bei den Hochschulen als auch bei Bund und Ländern, die die öffentlichen Hochschulen zu finanzieren haben.
